\chapter{Introduzione}

\section{Contesto}

Il \textbf{Registro Nazionale degli Aiuti} (RNA) è la banca dati nazionale istituita presso il Ministero dello Sviluppo Economico che raccoglie le informazioni relative agli aiuti di Stato e agli aiuti \textit{de minimis} concessi in Italia. L'RNA pubblica periodicamente dataset in formato \textbf{Open Data XML}, consentendo l'accesso pubblico a informazioni dettagliate su milioni di contributi erogati.

Questi dataset presentano caratteristiche che li rendono complessi da gestire:

\begin{itemize}
    \item \textbf{Volumi elevati}: centinaia di file XML, ciascuno contenente decine di migliaia di record
    \item \textbf{Struttura gerarchica}: dati annidati su più livelli (Aiuti $\rightarrow$ Componenti $\rightarrow$ Strumenti)
    \item \textbf{Qualità variabile}: presenza di caratteri XML non validi e encoding inconsistente
    \item \textbf{Assenza di formati analitici}: i dati sono distribuiti solo in XML, non ottimale per query analitiche
\end{itemize}

\section{Obiettivi del Progetto}

\textbf{Open Data Chunker} è una pipeline ETL (\textit{Extract, Transform, Load}) progettata per rispondere a queste sfide. Gli obiettivi principali sono:

\begin{enumerate}
    \item \textbf{Parsing robusto}: capacità di processare file XML corrotti o malformati senza interruzioni
    \item \textbf{Performance scalabili}: elaborazione parallela multi-processo per sfruttare CPU multi-core
    \item \textbf{Output analitico}: conversione in formato \textbf{Apache Parquet} con partizionamento temporale
    \item \textbf{Query interattive}: integrazione con motori SQL (\textbf{DuckDB}) per interrogazioni ad-hoc
    \item \textbf{Riproducibilità}: containerizzazione Docker per ambienti consistenti
\end{enumerate}

\begin{infobox}[Punti di Forza]
\begin{itemize}
    \item Pipeline end-to-end: dal raw XML al dataset queryable
    \item Fault-tolerant: gestione automatica degli errori
    \item Compressione 10x rispetto all'XML originale
    \item Deploy immediato via Docker
\end{itemize}
\end{infobox}

\section{Target Audience}

Questa documentazione è rivolta a \textbf{Software Engineer} e \textbf{Data Engineer} interessati a:

\begin{itemize}
    \item Comprendere l'architettura di una pipeline ETL Python moderna
    \item Valutare pattern di parallelismo e gestione memoria
    \item Replicare o estendere la soluzione per casi d'uso simili
\end{itemize}
